\section{Introduction}
\label{sec:introduction-old}

Session types
\cite{DBLP:conf/concur/Honda93,DBLP:conf/parle/TakeuchiHK94} provide a
rich type discipline for communication channels that guarantees
protocol fidelity statically. Many typing features have been
translated to the setting of session types, among them subtyping
\cite{DBLP:journals/acta/GayH05}, polymorphism
\cite{lindley17:_light_funct_session_types}, dependent types
\cite{DBLP:conf/ppdp/ToninhoCP11,DBLP:conf/fossacs/ToninhoY18,DBLP:journals/pacmpl/ThiemannV20},
and gradual types \cite{DBLP:journals/jfp/IgarashiTTVW19}. Recursive
types have been added at a very early stage
\cite{DBLP:conf/esop/HondaVK98} in a tail-recursive and equirecursive
fashion. Tail recursion was dictated by the asymmetric sequential manner of type
formation, inherited from the $\pi$-calculus. Equirecursion, where a
recursive type is considered equal to its unfoldings (or even to its
infinite unfolding), was likely chosen because unfolding should not
give rise to communication.  Moreover, there are efficient algorithms
for deciding equivalence and subtyping of recursive session types in
this form because these problems can be reduced to decision problems
on regular languages. 

However, there are protocols that cannot easily be described with
regular, tail-recursive structures. Examples are the serialization of
tree structures, stack-based protocols, non-regular protocols for
components \cite{DBLP:conf/soco/Sudholt05}, and typestate structures expressed
using fluent interfaces
\cite{DBLP:conf/ecoop/GilL16,DBLP:conf/oopsla/Xie17}.
Context-free and nested session types have been developed to  describe
these protocols
\cite{DBLP:conf/icfp/ThiemannV16,DBLP:conf/esop/Padovani17,DBLP:journals/toplas/Padovani19,DBLP:journals/corr/abs-2106-06658, DBLP:conf/esop/DasDMP21}. These
systems either equip the syntax of session types with a symmetric monoidal structure,
thus freeing them from the asymmetry of previous systems or allow
nesting of type constructors. Adding
equirecursion to such a system gives rise to context-free protocol
structures. Unfortunately, for context-free as well as nested
session types, this combination of features complicates type equivalence
significantly (doubly exponential \cite{DBLP:conf/esop/DasDMP21}) and
renders subtyping undecidable \cite{DBLP:journals/toplas/Padovani19}.

The following code example \cite{DBLP:conf/icfp/ThiemannV16} sends a
serialized binary tree using a context-free session type:
\begin{lstlisting}
data Ast = Con Int | Add Ast Ast
sendAst : Ast -> forall(s:S). mua. +{PCon: !Int, PAdd: a; a}; s -> s
sendAst t [s] c = case t of {
  Con x     -> select PCon [s] c
             |> send [Int, s] x,
  Add tl tr -> select PAdd [s] c
             |> sendAst tl [mua. +{PCon: !Int, PAdd: a; a};s]
             |> sendAst tr [s] }
\end{lstlisting}
The reading of \lstinline+data+ and \lstinline+case+ is inspired by
Haskell. The key for sending the serialized tree is the recursive session type
\lstinline_mua. +{PCon: !Int, PAdd: a; a}_. The choice operator $\oplus$
indicates a choice of sending one of the tags \lstinline+PCon+ and
\lstinline+PAdd+ and then continuing according to the session type
after the colon. For \lstinline+PCon+,
it sends the payload as determined by \lstinline+!Int+. For \lstinline+PAdd+, it recursively runs the
session type twice: \lstinline+a; a+. The semicolon operator
$\TSeq\_\_$  stands for sequential composition of session types. The operations on the
type are the usual ones: \lstinline+select+ chooses one of the tags and
sends it on the channel, \lstinline+send+ takes a number and sends it.
The operator \lstinline+|>+ is reversed function application and the
square brackets indicate type abstraction and type application.

Type checking this code in the system of \citet{DBLP:conf/icfp/ThiemannV16} is
already quite demanding. It requires unfolding of the recursion,
associativity of the $\TSeq\_\_$ operator, and distributivity of
$\TSeq\_\_$ over the choice operator. Moreover, the type checker
incorporates an expensive algorithm for type equivalence. 

A similar situation with recursive types promoted the invention of algebraic datatypes
\cite{Burstall1977a,DBLP:conf/lfp/BurstallMS80}, in a time when the decision problem
underlying type equivalence (equivalence of deterministic pushdown systems) was still open.
Algebraic datatypes combine sum-of-product types with recursion in a
way that makes type equivalence easy. In fact, the \lstinline+Ast+
datatype is isomorphic to the equirecursive type
\lstinline_mua.Int + atimesa_ where the body is a sum of product
types. Type equivalence for recursive types like this one runs into
similar difficulties with type equivalence as context-free session
types. The contribution of algebraic datatypes is the ``shrink
wrapping'' of recursion with sum-of-product types in combination with
naming the type. This combination makes type checking for algebraic datatypes
nominal and thus sidesteps all difficulties arising from recursive
types. 

The innovation of algebraic protocols consists of introducing a
shrink wrapping of recursion with choice/branch-of-sequence types combined
with naming. Compared with algebraic datatypes, we move from sum to
choice/branch and from product to sequence, which fits precisely with
the linear logic interpretation of session types \cite{DBLP:journals/mscs/CairesPT16}. 
% They implement isorecursion by reusing the
% choice tags for folding and the branch tags for driving unfolding. 
As an example, we review the serialization of binary trees using an
algebraic protocol. Like an algebraic datatype, we name the protocol,
we tag the alternatives, and we define the payload as a sequence:
\begin{lstlisting}
protocol PAst = PCon Int | PAdd PAst PAst
\end{lstlisting}
This definition only introduces two alternatives and fixes the
sequencing, but it does not indicate a direction of communication,
yet. To this end, we reuse the standard session type notation:
\begin{itemize}
\item \lstinline/!PAst/ declares sending: send a tag followed by the payload;
\item \lstinline/?PAst/ declares receiving: receive a tag, branch
  accordingly, and receive the payload.
\end{itemize}
Now we can revisit the function to send the binary tree:
\begin{lstlisting}
sendAst : Ast -> forall (s:S). !PAst.s -> s
sendAst t [s] c = case t of {
  Con x     -> select PCon [s] c
             |> send [Int, s] x,
  Add tl tr -> select PAdd [s] c
             |> sendAst tl [!PAst.s] |> sendAst tr [s] }
\end{lstlisting}
Perhaps shockingly, the code remains the same. But the types change in
several respects.
\begin{itemize}
\item The symmetric sequencing operator $\TSeq\_\_$ disappears from the syntax along with all
  problems posed by its associativity and distributivity. The
  \emph{syntax} of session types is tail-recursive as in $\TIn TS$,
  again. The context-free behavior is achieved by the typing of
  \lstinline/select PAdd/. Its type \lstinline/forall(s:S).!PAst.s -> !PAst.!PAst.s/
  ``rewrites'' a single occurrence of \lstinline/!PAst/ in front of
  the session type to two occurrences corresponding to the payload of
  the tag \lstinline/PAdd/, like a pushdown automaton operating on
  the session type.
\item Checking equivalence of algebraic protocol types is efficient:
  it is achieved by comparing their names.
\item The type \lstinline/!PAst.s/ implies that after selecting a tag,
  the payload is also sent. This mode extends to the recursive uses of
  \lstinline/PAst/ in the protocol. We would use \lstinline/?PAst.s/
  to receive a binary tree. 
\end{itemize}
This example gives a taste of the power of algebraic protocols.
In \cref{sec:algebr-sess-types} we explain how to change direction in
a protocol and how to write modular protocols by exploiting parametric
protocols.



\subsection*{Contributions}
\label{sec:contributions}

\begin{itemize}
\item Algebraic protocols (AP). We propose a new approach for defining
  recursive protocols compositionally inspired by parameterized algebraic
  datatypes. This approach combines naming, recursion, and
  choice/branch-of-sequence types analogously to algebraic
  datatypes. It adds a novel notion of direction to the payload of a
  tag. Checking type equivalence becomes simple and efficient
  (linear in the size of the types) as APs are nominal. 
\item Algebraic session types (\algst). We define a (mostly) linearly typed core calculus that
  is based on linear System~F with recursion and features session
  types based on APs. Type checking for \algst is specified  using bidirectional
  typing rules. 

  The operational semantics is presented as a labelled transition relation, adapted
  from Fowler \etal and Montesi and
  Peressotti~\cite{DBLP:conf/concur/0001KDLM21,DBLP:journals/corr/abs-2105-08996,DBLP:journals/corr/abs-2106-11818}
  for double binders and the open/close scheme for free output of the
  $\pi$-calculus~\cite{DBLP:journals/iandc/MilnerPW92b}.
  We prove type soundness via preservation and progress, where our approach offers a
  rather elegant statement (\cref{thm:progress} on \cpageref{thm:progress}). 
\item Implementation. The proposed system is fully implemented
  including several practical extensions that are documented in
  \cref{sec:implementation}. All examples in
  \cref{sec:algebr-sess-types} are checked on the implementation.
\end{itemize}

\subsection*{Overview}
\label{sec:overview}

\Cref{sec:algebr-sess-types} contains an informal introduction of
algebraic protocols with examples. \Cref{sec:types} formally
introduces kinds and types of \algst and \cref{sec:processes} defines
the statics and dynamics of expression and
processes. \Cref{sec:metatheory} presents the metatheoretical
results.
% \Cref{sec:cfst-vs-algst} gives an informal outline of
% an embedding of context-free session types to \algst.
\Cref{sec:implementation} discusses additional features of the
implementation. \Cref{sec:related,sec:conclusion} discuss related
work and conclude.

The proofs and auxiliary results are provided as supplementary material. 
The implementation is available for review in an anonymous repository
and will be submitted for artifact evaluation:
\begin{center}
  \url{https://gitlab.com/algst/algst}
\end{center}

%%% Local Variables:
%%% mode: latex
%%% TeX-master: "main"
%%% End:
